% ============================================================================
%  Мастер рад — Никола Суботић
%  Решавање једнодимензионог проблема паковања применом генетског алгоритма
%
%  Прављење (xelatex није инсталиран, користи се pdflatex + bibtex):
%     build.bat master-thesis
%  или ручно:
%     pdflatex master-thesis  ->  bibtex master-thesis  ->  pdflatex x2
% ============================================================================
\documentclass[12pt,oneside]{memoir}

% Стил мастер рада Математичког факултета (подразумевано: ћирилица, bibtex)
\usepackage{matfmaster}

% Исправан приказ ћириличких италик слова под pdflatex-ом
\usepackage{cmsrb}

% Математички пакети
\usepackage{amsmath}
\usepackage{amssymb}

% Датотека са литературом (master-thesis.bib)
\bib{master-thesis}

\autor{Никола Суботић}
\naslov{Решавање jеднодимензионалног проблема паковања у кутиjе коришћењем генетског алгоритма}
\godina{2026}
\mentor{проф. др Мирослав \textsc{Марић}\\ Универзитет у Београду, Математички факултет}
\komisijaA{др Нина \textsc{Радојичић Матић}\\ Универзитет у Београду, Математички факултет}
\komisijaB{Денис \textsc{Аличић}\\ Универзитет у Београду, Математички факултет}
% \datumodbrane{...}

% Апстракт (по MSNR правилу пише се последњи — ово је радна верзија)
\apstr{%
Овај рад разматра једнодимензиони проблем паковања, NP-тежак проблем комбинаторне
оптимизације. Предложен је приступ заснован на генетском алгоритму са пермутационим
кодирањем и \foreignlanguage{english}{First-Fit} декодером, чија је почетна популација
засејана класичним хеуристикама. Приступ је упоређен са класичним хеуристикама
(\foreignlanguage{english}{First-Fit}, \foreignlanguage{english}{Best-Fit},
\foreignlanguage{english}{Next-Fit}, \foreignlanguage{english}{Max-Rest} и њиховим
варијантама са сортирањем) на стандардним скуповима инстанци.
}

\kljucnereci{једнодимензиони проблем паковања, генетски алгоритам, хеуристике,
First-Fit, комбинаторна оптимизација}

% ============================================================================
\begin{document}

% ---- Уводни део ------------------------------------------------------------
\frontmatter
\naslovna
\komisija
% \posveta{...}
\apstrakt
\tableofcontents*

% ---- Главни део ------------------------------------------------------------
\mainmatter

% Увод — ненумерисано поглавље (по MSNR пише се последње; ово је радна верзија).
% Роадмап реферише на поглавља 2--6 која још нису написана; ускладити при финализацији.
\chapter*{Увод}
\markboth{УВОД}{УВОД}
\phantomsection
\addcontentsline{toc}{chapter}{Увод}

Проблем паковања представља један од класичних проблема комбинаторне оптимизације: дат
је скуп предмета познатих величина које треба распоредити у што мањи број кутија
ограниченог капацитета. Иако се лако формулише, овај проблем се појављује у бројним
практичним ситуацијама --- од сечења материјала стандардних димензија и утовара робе у
возила ограничене носивости, до распоређивања послова на процесорима и доделе меморије у
рачунарским системима. У свим тим ситуацијама свака уштеђена кутија значи мање утрошеног
материјала, простора или времена, па је проналажење што бољег распореда од непосредног
практичног значаја.

Једнодимензиони проблем паковања припада класи NP-тешких проблема, због чега се за
инстанце реалних димензија не може очекивати да ће егзактне методе пронаћи оптимално
решење у прихватљивом времену. Из тог разлога се у пракси примењују хеуристике, које брзо
дају добра, али не нужно оптимална решења, и метахеуристике, које усмереном претрагом
простора решења теже да таква решења додатно побољшају.

Предмет овог рада јесте решавање једнодимензионог проблема паковања генетским алгоритмом.
Предложени алгоритам користи пермутационо кодирање, у коме хромозом представља редослед
предмета, док се конкретно паковање добија применом хеуристике First-Fit на тај редослед.
Почетна популација засејава се решењима класичних хеуристика, чиме се, уз елитизам,
обезбеђује да генетски алгоритам никада не да лошије решење од најбоље хеуристике којом је
засејан. Предложени приступ упоређује се са класичним хеуристикама --- First-Fit,
Best-Fit, Next-Fit и Max-Rest, као и њиховим варијантама са претходним сортирањем --- на
стандардним скуповима инстанци из библиотеке BPPLIB~\cite{delorme2018}.

Рад је организован на следећи начин. У првом поглављу уводи се формална дефиниција
једнодимензионог проблема паковања, даје се преглед досадашњих приступа његовом решавању
и наводе се области примене. Друго поглавље описује класичне хеуристике и генетски
алгоритам за груписање као полазиште за метахеуристички приступ. У трећем поглављу излаже
се предложени хибридни генетски алгоритам са пермутационим кодирањем и хеуристичким
засејавањем почетне популације. Четврто поглавље посвећено је детаљима имплементације, а
пето поглавље приказује и анализира експерименталне резултате добијене поређењем са
класичним хеуристикама и познатим оптималним решењима. У закључку се сумирају резултати
рада и наводе правци даљег истраживања.


\input{poglavlje1}

% TODO: поглавља 2--6 (хеуристике, методологија, имплементација, резултати, закључак)

% ---- Литература ------------------------------------------------------------
\literatura

% ---- Завршни део -----------------------------------------------------------
\backmatter
\begin{biografija}
% TODO: биографија аутора
\noindent Овде долази кратка биографија аутора.
\end{biografija}

\end{document}
